\section{Introduction}

The use of Artificial Intelligence (AI) into healthcare environment
presents a exceptional opportunities to optimize processes, improve
the diagnosis precision and leverage healthcare staff from executing
repetivive tasks to provide a better treat to patients. Unlike other
machine learning applications, medical data cannot be easily
centralized for training models, as strict privacy policies are
imposed to protect the patient. Furthermore, clinical information is
inherently distributed across multiple distinct hospitals and
organizations, and the datasets themselves are frequently unbalanced
due to the scarcity of specific diagnostic classes or rare diseases.

This mini-project illustrates the core challenges and proposes a
modular, end-to-end AI pipeline designed to overcome them. To validate
the architecture, we selected the binary classification of Chest
X-Rays (Pneumonia vs. Normal). While modern Convolutional Neural
Networks (CNNs) can easily achieve high accuracy on this specific
task, the dataset presents a significant class imbalance due to a
scarcity of the minority 'NORMAL' class, making it an ideal candidate
for our project.

The true goal of this report is to demonstrate a framework that
navigates privacy, distribution, and data scarcity. To achieve this,
the proposed system architecture leverages three distinct
techniques: 
\begin{itemize}
\item \textbf{Generative Adversarial Networks (GAN):}
\item \textbf{Convolutional Neural Networks (CNN):}
\item \textbf{Federated Learning (FL)}: 
\end{itemize}